\chapter[Discussão e Conclusão]{Discussão e Conclusão}
\label{cap:conclusao}

Este capítulo retoma os resultados apresentados no Capítulo~\ref{cap:resultados} e os relaciona ao problema de pesquisa formulado na introdução. O objetivo central era verificar em que medida uma arquitetura supervisionada especializada, baseada em backbone facial, descongelamento parcial e tokens regionais anatômicos, apresenta ganho em relação a uma linha de base sem treinamento e a outras arquiteturas supervisionadas no FIW. A comparação foi feita sob um protocolo único de verificação binária. Além disso, o trabalho avaliou de forma exploratória modelos multimodais generalistas como linhas de base binárias zero-shot, mantendo a avaliação multiclasse anterior apenas como diagnóstico histórico.

\section{Discussão dos achados}

O principal achado experimental é que o AdaFace-Regional com Negativos Difíceis alcançou o maior AUC global entre as abordagens avaliadas sob validação cruzada. Essa configuração obteve AUC de 0,876 $\pm$ 0,003, acima do ViT-FaCoR (0,846 $\pm$ 0,004), do DINOv2-LoRA com descongelamento completo (0,814 $\pm$ 0,007), da linha de base AdaFace-Cosseno (0,799 $\pm$ 0,000) e do Modelo com Recuperação (0,778 $\pm$ 0,008). A diferença em relação ao AdaFace-Cosseno é importante porque essa linha de base já usa um extrator facial forte, treinado para reconhecimento de identidade. Portanto, o ganho do AdaFace-Regional está associado ao conjunto formado pela adaptação desse backbone à tarefa de parentesco, comparação regional e passagem simétrica.

Em um protocolo externo zero-shot, sem validação cruzada, o Claude Sonnet binário alcançou AUC de 0,789 e o Codex VLM binário alcançou AUC de 0,624. Esses resultados ajudam a situar os VLMs, mas não substituem a comparação principal entre modelos avaliados sob validação cruzada.

A comparação com o ViT-FaCoR também responde ao eixo principal do trabalho. O ViT-FaCoR representa a referência supervisionada mais forte entre as arquiteturas anteriores avaliadas e fica cerca de 0,030 abaixo no valor de AUC em relação à configuração com Negativos Difíceis. Essa diferença é maior do que a variação observada entre as dobras, repete-se nas configurações Simétrica e Comparativa do AdaFace-Regional e é significativa no teste t pareado sobre AUC para a configuração com Negativos Difíceis ($t(4)=10{,}264$, $p=0{,}0005$). No Wilcoxon bilateral, entretanto, a comparação fica em $p=0{,}0625$ com apenas cinco dobras. Assim, a evidência experimental indica uma vantagem empírica da proposta autoral na ordenação global dos pares no FIW dentro do protocolo adotado.

Ao mesmo tempo, os resultados não indicam uma superioridade simples de uma única configuração em todos os regimes. As configurações Simétrica, Comparativa e com Negativos Difíceis ficam muito próximas em AUC, e a escolha entre elas depende mais do ponto de operação do que de uma diferença global clara. Entre as variantes do AdaFace-Regional, a configuração com Negativos Difíceis é mais forte em TAR@FAR=0{,}001, enquanto a Comparativa e a com Negativos Difíceis deslocam o modelo para regimes de menor falsa aceitação em comparação com a Simétrica. Esse ponto é relevante porque, em aplicações forenses ou humanitárias, falsos positivos podem ter custo maior do que falsos negativos. Nesses casos, AUC alto não substitui a escolha explícita de um limiar e de uma taxa de falsa aceitação aceitável.

No ciclo do AdaFace-Regional, a mudança mais clara foi a passagem simétrica. O modelo com descongelamento parcial já chegava a 0,856 de AUC, mas ainda tinha um problema visível nas relações de avós e netos. Quando o modelo passou a treinar o mesmo par nas duas ordens, $(A,B)$ e $(B,A)$, o AUC médio em validação cruzada subiu para 0,873 $\pm$ 0,004, e essas relações de segundo grau ficaram entre 80,0\% e 84,3\% de acurácia. A interpretação mais plausível é que o aprendizado do modelo anterior era enviesado por artefatos estatísticos ligados à posição das imagens no par. A simetrização reduz esse atalho ao tornar o escore final menos dependente da ordem de entrada.

Também é importante analisar o AdaFace-Cosseno com cuidado. Ele não foi treinado para parentesco, mas ainda chega a 0,218 $\pm$ 0,000 em TAR@FAR=0{,}01 e a 0,071 $\pm$ 0,000 em TAR@FAR=0{,}001. Ou seja, quando o limiar é mais conservador, o embedding de identidade do AdaFace já separa parte dos pares. O AdaFace-Regional parte desse sinal e melhora a ordenação geral, além de recuperar melhor classes positivas difíceis, principalmente avós e netos.

Nos VLMs, a leitura ficou mais informativa após a reformulação binária. O Claude Sonnet, em zero-shot, chega a AUC de 0,789 e fica próximo da linha de base AdaFace-Cosseno, mas ainda abaixo dela. O Codex VLM fica bem mais baixo, com AUC de 0,624, apesar de recuperar muitos pares positivos. O padrão de erro mais importante é a dificuldade de rejeitar non-kin. O Codex aceita como parentes 1.857 dos 3.000 pares negativos, e o Claude ainda produz 1.082 falsos positivos. Além disso, ambos têm TAR@FAR=0{,}01 baixo, o que mostra que a confiança textual do VLM não gera uma ordenação suficientemente fina para pontos operacionais conservadores.

Essa leitura é reforçada pela avaliação adicional do AdaFace-Regional R012 no mesmo manifesto de 6.000 pares usado pelo Claude Sonnet. Nesse recorte controlado, o modelo supervisionado alcança AUC de 0,888, Average Precision de 0,882 e TAR@FAR=0{,}01 de 0,259, contra 0,789, 0,742 e 0,042 do Claude. Esse resultado não substitui a validação cruzada principal, mas mostra que a diferença em relação ao VLM permanece quando ambos são avaliados sobre o mesmo conjunto de pares. O ensaio multiclasse descrito na metodologia reforça essa leitura como diagnóstico auxiliar. Os resultados indicaram baixa capacidade de distinguir relações positivas finas, especialmente no recorte de avós e netos, e por isso esse ensaio não compõe a comparação principal.

\section{Contribuições do trabalho}

A primeira contribuição do trabalho é o protocolo comparativo uniforme. As cinco abordagens de verificação binária foram avaliadas no FIW sob a mesma lógica de validação cruzada por família, seleção de limiar em validação e métricas independentes de limiar. Essa uniformização reduz uma fragilidade frequente na literatura de parentesco facial, em que resultados obtidos com partições, limiares e métricas diferentes acabam sendo comparados diretamente.

A segunda contribuição é a proposta arquitetônica do AdaFace-Regional. O modelo combina AdaFace IR-101 parcialmente descongelado, cinco tokens regionais anatômicos, atenção cruzada entre regiões, gate sigmoide, cabeça auxiliar de relação e passagem simétrica. A proposta não é apenas trocar o backbone por outro mais forte. Ela introduz uma forma explícita de comparar regiões faciais e de reduzir a dependência da decisão em relação às duas ordens do par, propriedade esperada em uma tarefa de verificação de parentesco.

A terceira contribuição é a análise por tipo de relação. Em vez de discutir apenas métricas globais, o trabalho mostra como o desempenho muda entre pais e filhos, mães e filhos, irmãos, avós e netos e a classe non-kin. Essa análise é importante porque o FIW não tem dificuldade uniforme entre relações. As classes de avós e netos são menores e historicamente mais instáveis, e justamente nelas o AdaFace-Regional apresenta um dos ganhos mais relevantes.

A quarta contribuição é a avaliação exploratória dos VLMs em formulação binária. O trabalho registra como Codex e Claude Sonnet se comportam quando recebem a mesma pergunta kin/non-kin dos modelos supervisionados, mas sem treino, validação cruzada ou calibração de limiar. Esse recorte ajuda a mostrar o limite da inferência direta por prompt. Os modelos capturam parte do sinal visual, mas ainda não dão evidência de que possam substituir um modelo supervisionado treinado para parentesco facial.

\section{Limitações}

Os resultados devem ser interpretados dentro do escopo experimental definido. A validação cruzada de cinco dobras em FIW reduz a dependência de uma única partição, mas não elimina todas as fontes de variação. Os testes estatísticos pareados foram calculados sobre apenas cinco valores por modelo, um por dobra, e, portanto, têm baixo poder amostral; no Wilcoxon bilateral, por exemplo, o menor p-valor exato possível com cinco dobras é 0,0625. A principal limitação remanescente é a ausência de repetição por múltiplas sementes em cada dobra. Por isso, diferenças pequenas entre as configurações Simétrica, Comparativa e com Negativos Difíceis devem ser lidas junto com o ponto operacional, não como superioridade absoluta de uma configuração sobre as outras.

Outra limitação é a generalização externa. O trabalho não avalia transferência para bases fora do FIW, imagens de baixa qualidade, populações demograficamente distintas ou cenários operacionais reais. A literatura de reconhecimento facial mostra que esse tipo de mudança pode alterar bastante o desempenho, e o mesmo risco vale para parentesco facial. Também não foi feita uma auditoria de sobreposição entre o WebFace4M, usado no pré-treino do AdaFace, e identidades ou imagens presentes no FIW. Não há evidência local de vazamento, mas a possibilidade permanece como risco metodológico não verificado.

Devido à integração simultânea de múltiplos componentes metodológicos, o desenho experimental atual carece de um estudo de ablação estrito, o que impede a quantificação da contribuição isolada de cada módulo, como atenção regional, gate sigmoide, cabeça auxiliar de relação e passagem simétrica. O resultado mais seguro é dizer que o conjunto formado por descongelamento parcial, regiões faciais, cabeça auxiliar e passagem simétrica está associado ao ganho do AdaFace-Regional no protocolo usado. Trabalhos futuros devem desmembrar essas implementações para validar seus efeitos de forma independente.

Uma limitação específica está no histórico de amostragem de negativos. O ciclo experimental incluiu execuções preliminares de reamostragem que não foram usadas para sustentar a conclusão sobre negativos difíceis. A análise final do AdaFace-Regional considera como negativos difíceis somente a configuração posterior com pares negativos reais, formados por pessoas de famílias diferentes com os mesmos papéis do par positivo.

A avaliação com VLMs tem escopo mais restrito do que a comparação supervisionada. Embora a tarefa final seja binária, os VLMs foram avaliados em uma única execução sobre 6.000 pares balanceados, sem validação cruzada, sem ajuste de limiar e sem repetição temporal. A execução adicional do AdaFace-Regional R012 no manifesto usado pelo Claude controla o conjunto de pares para essa comparação específica, mas também não substitui o teste oficial completo nem a validação cruzada. Além disso, os modelos avaliados são proprietários e podem mudar com atualizações de fornecedor, o que reduz a estabilidade dos resultados. A avaliação multiclasse anterior permanece ainda mais limitada. Ela não inclui non-kin, usa amostras menores e serve apenas como diagnóstico de classificação fina de relações.

\section{Implicações práticas e éticas}

A verificação visual de parentesco facial pode ser útil como ferramenta auxiliar em triagem, busca familiar e apoio investigativo. No entanto, os resultados deste trabalho não autorizam uso decisório isolado. Mesmo o melhor modelo ainda produz falsos positivos e falsos negativos, e esses erros têm custo alto em contextos forenses e humanitários, especialmente quando envolvem populações vulneráveis.

Fora do ambiente experimental, a saída do modelo deve ser lida apenas como um indício de triagem. A decisão final precisa caber a uma equipe qualificada e deve ser confrontada com DNA, documentação oficial, investigação social ou avaliação especializada. Como as imagens de face são dados biométricos, qualquer aplicação real também teria que prever controle de acesso, registro do uso, possibilidade de contestação, consentimento quando aplicável e base legal clara. Antes de uso em população real, seria necessário acompanhar os erros por grupo demográfico, porque diferenças de desempenho nesse tipo de sistema podem atingir justamente os grupos mais vulneráveis.

\section{Conclusão}

O objetivo geral do trabalho foi comparar, sob um protocolo comum, uma proposta supervisionada especializada para verificação visual de parentesco facial com linhas de base e arquiteturas alternativas. Esse objetivo foi atendido dentro do escopo definido. No FIW, o AdaFace-Regional com Negativos Difíceis alcançou AUC de 0,876 $\pm$ 0,003 em validação cruzada, ficando cerca de 0,030 acima do ViT-FaCoR no valor de AUC, com suporte do teste t pareado ($p=0{,}0005$), embora o Wilcoxon bilateral não atinja 0,05 com cinco dobras. A proposta também fica 0,062 acima do DINOv2-LoRA, 0,077 acima da linha de base AdaFace-Cosseno e 0,098 acima do Modelo com Recuperação. As linhas de base VLM binárias zero-shot também ficaram abaixo da proposta, mas devem ser lidas como avaliação externa em outro protocolo.

Os resultados sustentam, dentro do protocolo adotado, ganho em AUC global para o AdaFace-Regional. A arquitetura proposta fica acima da linha de base sem treinamento e das demais arquiteturas supervisionadas avaliadas quanto ao AUC global. Na análise por relação, o aspecto mais relevante é a recuperação das classes de avós e netos nas variantes mais voltadas às classes positivas, justamente o grupo que aparecia como mais frágil nas primeiras execuções do modelo.

Em síntese, embora o resultado final aponte avanços relevantes, também impõe uma ressalva importante. Um AUC mais alto, isoladamente, não resolve a escolha do ponto de operação. O AdaFace-Cosseno ainda apresenta bom desempenho em regimes de baixa falsa aceitação, enquanto as versões do AdaFace-Regional alteram seu comportamento em função do limiar e do custo dos erros.

Os VLMs binários ajudam a situar o quanto a tarefa pode ser respondida por inferência multimodal genérica, mas seus erros em non-kin e seus baixos valores de TAR@FAR=0{,}01 reforçam a necessidade de modelos especializados. A comparação adicional no manifesto usado pelo Claude confirma essa leitura em um subconjunto controlado de pares, no qual o AdaFace-Regional R012 supera o VLM em todas as métricas reportadas. Dessa forma, as principais contribuições do trabalho são a comparação controlada entre abordagens, a proposição de uma arquitetura regional associada a melhor desempenho no FIW, a avaliação binária de VLMs zero-shot e uma análise segmentada por tipo de relação, destacando tanto os avanços observados quanto as limitações que permanecem na tarefa.
